\begin{abstract}
Recent advancements, exemplified by BitNet~\cite{bitnet}, signal the emergence of a new generation of 1-bit Large Language Models (LLMs). Here, we propose \textbf{\text{BitNet b1.58}}, a 1-bit LLM variant wherein every individual parameter (or weight) adopts a ternary value from the set \{-1, 0, 1\}. This model achieves parity with full-precision (e.g., FP16 or BF16) Transformer LLMs of equivalent scale and exposed to the same volume of training data, matching their perplexity and downstream task results, yet at a fraction of the cost in latency, memory consumption, throughput, and energy use. Notably, the 1.58-bit LLM introduces a revised scaling law and a novel methodology for training next-generation LLMs that balance state-of-the-art performance with efficiency. In addition, it catalyzes the adoption of a new computational paradigm and paves the way for the development of bespoke hardware tailored for 1-bit LLM operations.
\end{abstract}

\section{The Era of 1-bit LLMs}

The artificial intelligence community has recently witnessed explosive advancements in both the magnitude and capabilities of Large Language Models (LLMs). While these models set new benchmarks across an array of natural language processing applications, their expanding scale has introduced formidable barriers to practical deployment, as well as escalating apprehensions regarding environmental and economic costs tied to significant energy demands.

One promising avenue to alleviate such obstacles involves the application of post-training quantization strategies to realize low-bit inference models~\cite{smoothquant, gptq, quip, quip_sharp}. By lowering the numerical precision of model weights and activations, this method substantially diminishes the memory footprint and computational burden associated with LLMs. The industry has notably shifted from the conventional 16-bit implementations towards even more compact representations, such as 4-bit configurations~\cite{gptq, awq}. Despite its broad adoption in industry LLMs, post-training quantization remains an imperfect solution.

Advancements in 1-bit model architectures, exemplified by BitNet~\cite{bitnet}, highlight a compelling avenue for decreasing the computational burden of LLMs while preserving their accuracy. Conventional LLMs utilize 16-bit floating-point formats (e.g., FP16 or BF16), where matrix multiplication operations dominate both the structure and computational demands of these models. As a result, a substantial proportion of resource expenditure is devoted to floating-point arithmetic. BitNet, on the other hand, conducts matrix multiplication through integer additions exclusively, drastically lowering the energy requirements compared to traditional floating-point operations. Since power often represents the primary constraint on chip performance, such energy reductions translate into accelerated computation.

Beyond computation, another significant bottleneck arises from moving model parameters from DRAM to accelerator-side memory (such as SRAM) during inference. Scaling up SRAM to alleviate this cost tends to be substantially more expensive than relying on DRAM. With their reduced memory footprint, 1-bit LLMs alleviate both capacity and bandwidth pressures relative to full-precision counterparts, resulting in more efficient and expedient weight loading, thus optimizing inference speed and efficiency.

This paper introduces an advanced 1-bit LLM variant, termed \textbf{\text{BitNet b1.58}}, where each parameter may assume one of three values: \{-1, 0, 1\}. By extending the original BitNet's 1-bit structure with the explicit inclusion of zero, this architecture corresponds to 1.58 bits in a binary system. \text{BitNet b1.58} preserves the foundational computational efficiencies of its 1-bit predecessor—including near-elimination of multiplication operations during matrix multiplications and associated opportunities for hardware-level optimization—while also retaining the significant energy savings and improvements in memory, throughput, and latency relative to FP16 LLMs. Notably, \text{BitNet b1.58} brings two new advantages: first, its expressiveness is heightened by directly supporting feature selection via weight-level zeroing, which facilitates meaningful performance enhancements for 1-bit LLMs. Second, our empirical results demonstrate that, for models as large as 3B parameters, \text{BitNet b1.58} can rival the perplexity and downstream performance of full-precision (FP16) baselines when controlling for configuration factors such as model size and amount of training data.

\section{BitNet b1.58}

\text{BitNet b1.58} builds upon the BitNet architecture, a variant of the Transformer in which the standard \emph{nn.Linear} modules are substituted by \emph{BitLinear}. The model is initialized with random parameters and trained using 1.58-bit quantized weights together with 8-bit quantized activations. Several distinctions from the original BitNet have been introduced and are outlined below.

\paragraph{Quantization Function.} To enforce weight values restricted to the set $\{-1, 0, +1\}$, an \emph{absmean} quantization function is implemented. This approach involves normalizing the weight matrix by its mean absolute value, followed by rounding each element to the closest integer among the permitted set:

\begin{equation}
    \widetilde{W} = \mathrm{RoundClip}(\frac{W}{\gamma + \epsilon}, -1, 1),
\end{equation}

\begin{equation}
    \mathrm{RoundClip}(x, a, b) = \max(a, \min(b, \mathrm{round}(x))),
\end{equation}

\begin{equation}
    \gamma = \frac{1}{nm}\sum_{ij} |W_{ij}|.
\end{equation}

The quantization method for activations remains consistent with BitNet, with the modification that, prior to applying non-linearities, activations are not normalized to the interval $[0, Q_b]$. Instead, activations are normalized to the range $[-Q_b, Q_b]$ on a per-token basis, effectively eliminating the need for zero-point quantization. This adjustment streamlines both the implementation and system-level optimizations, with empirical evidence showing negligible impact on performance.

\paragraph{LLaMA-alike Components.}

The LLaMA~\cite{llama, llama2} architecture has become the standard framework for open-source LLM research. To better align with open-source initiatives, the design of \text{BitNet b1.58} adopts LLaMA-like elements. This includes integrating RMSNorm~\cite{rmsnorm}, SwiGLU~\cite{swiglu}, rotary embeddings~\cite{rope}, and omitting all bias terms. Consequently, \text{BitNet b1.58} can be readily deployed within mainstream open-source platforms (such as Huggingface, vLLM~\cite{vllm}, and llama.cpp\footnote{\url{https://github.com/ggerganov/llama.cpp}}) with minimal modification. 

\section{Results}

We conducted a comparison between \text{BitNet b1.58} and our replicated FP16 \text{LLaMA LLM} models across several parameter scales. Both models underwent pre-training on the RedPajama dataset~\cite{redpajama} for a total of 100 billion tokens to guarantee equivalence in training conditions. For evaluation, we measured zero-shot performance over a diverse suite of language understanding tasks, such as ARC-Easy~\cite{arc}, ARC-Challenge~\cite{arc}, Hellaswag~\cite{hellaswag}, Winogrande~\cite{winoGrande}, PIQA~\cite{piqa}, OpenbookQA~\cite{openbookqa}, and BoolQ~\cite{boolq}. Additionally, we calculated validation perplexity on WikiText2~\cite{wikitext2} and C4~\cite{c4}.

We assessed the GPU memory consumption and inference latency for both \text{LLaMA LLM} and \text{BitNet b1.58}. Measurements were performed utilizing the FasterTransformer\footnote{\url{https://github.com/NVIDIA/FasterTransformer}} framework, which offers optimized inference for large language models on GPUs. For \text{BitNet b1.58}, we also integrated Ladder’s 2-bit kernel~\cite{ladder}. Our primary latency metric was the time required per generated output token, representing the predominant inference overhead.

A summary of perplexity and resource usage results for both \text{BitNet b1.58} and \text{LLaMA LLM} is provided in Table~\ref{tab:ppl}. These findings indicate that, beginning at the 3B model scale, \text{BitNet b1.58} achieves perplexity on par with its full-precision \text{LLaMA LLM} counterpart while delivering a 2.71-fold reduction in latency and consuming 3.55 times less GPU memory. Notably, the 3.9B variant of \text{BitNet b1.58} exhibits a 2.4$\times$ speedup, 3.32$\times$ lower memory usage, and a substantial accuracy improvement compared to \text{LLaMA LLM} 3B.

Comprehensive zero-shot accuracy outcomes for end tasks are detailed in Table~\ref{tab:zero-shot}. We employed the standardized \emph{lm-evaluation-harness}\footnote{\url{https://github.com/EleutherAI/lm-evaluation-harness}} procedure for assessment. Results reveal that as the scale increases, performance differences between \text{BitNet b1.58} and \text{LLaMA LLM} decrease. Crucially, from the 3B scale upward, \text{BitNet b1.58} consistently matches the full-precision baseline. Mirroring the perplexity findings, \text{BitNet b1.58} 3.9B surpasses the \text{LLaMA LLM} 3B model in final task performance, while maintaining advantages in memory and latency. These results position \text{BitNet b1.58} as a clear Pareto improvement relative to contemporary large language models.

\paragraph{Memory and Latency}

We extended our experiments to larger model sizes, specifically 7B, 13B, and 70B, and analyzed the associated computational costs. As depicted in Figure~\ref{fig:latency-memory-energy}, both latency and memory usage display clear scaling behavior: speed improvements become more pronounced as the number of parameters increases. Notably, \text{BitNet b1.58} at 70B parameters achieves a 4.1-fold reduction in latency relative to the \text{LLaMA LLM} reference. This improvement arises because the computational overhead for \emph{nn.Linear} operations grows proportionally with model scale. Memory usage patterns are comparable, primarily because embeddings use full-precision formats, but constitute a diminishing portion of total memory requirements as models get larger. The latency and memory benchmarks were performed with a 2-bit kernel, indicating there remain opportunities for further cost reduction with additional optimization.

\paragraph{Energy}

We additionally performed estimations of the energy consumption related to arithmetic operations for both \text{BitNet b1.58} and \text{LLaMA LLM}, with emphasis on matrix multiplication, given its dominant contribution to the operational cost of large language models. Figure~\ref{fig:energy} breaks down the sources of energy use. Most of the energy expenditure in \text{BitNet b1.58} arises from INT8 addition operations, in contrast to the mixed FP16 addition and multiplication in \text{LLaMA LLM}. Drawing on the energy estimation framework outlined in~\cite{energycost, pokebnn}, we find that \text{BitNet b1.58} reduces the energy required for matrix multiplication arithmetic operations by a factor of 71.4 on 7nm hardware. End-to-end energy assessments, measured with 512 token sequences, reveal that the efficiency advantage of \text{BitNet b1.58} grows with model size when compared to the FP16 \text{LLaMA LLM}. This increasing efficiency results from the greater share of \emph{nn.Linear} operations in larger models, while the relative computational demands of other model components decline.

\paragraph{Throughput}

We assess the throughput performance of \text{BitNet b1.58} and \text{LLaMA LLM}—each with 70B parameters—across two 80GB A100 GPUs, leveraging pipeline parallelism~\cite{gpipe} to accommodate the large-scale \text{LLaMA LLM} 70B model. The batch size was increased incrementally until the GPU memory reached its capacity, maintaining a fixed sequence length of 512. As illustrated in Table~\ref{tab:throughput}, \text{BitNet b1.58} 70B can sustain a batch size up to 11 times larger than \text{LLaMA LLM} 70B, leading to an 8.9-fold improvement in throughput.

\textbf{\text{BitNet b1.58} introduces a new paradigm regarding the tradeoff between model effectiveness and inference expense.} According to Figure~\ref{fig:latency-memory-energy} and Figure~\ref{fig:energy}, we can summarize equivalences between 1.58-bit and 16-bit models as follows:

\begin{itemize}
    \item 13B BitNet b1.58 demonstrates greater efficiency—including lower latency, memory requirements, and energy usage—than a 3B FP16 LLM.
    \item 30B BitNet b1.58 outperforms a 7B FP16 LLM in terms of latency, memory consumption, and energy cost.
    \item 70B BitNet b1.58 surpasses a 13B FP16 LLM across latency, memory usage, and energy efficiency metrics.
\end{itemize}

\paragraph{Training with 2T Tokens}

For large language models, the volume of training tokens is a significant determinant of model capability. To evaluate how \text{BitNet b1.58} scales with token count, we trained the model using 2T tokens following the same dataset construction strategy as StableLM-3B~\cite{StableLM-3B-4E1T}, which represents the leading open-source 3B LLM. Both models were assessed using a benchmark suite that includes Winogrande~\cite{winoGrande}, PIQA~\cite{piqa}, SciQ~\cite{sciq}, LAMBADA~\cite{lambada}, and ARC-easy~\cite{arc}. Zero-shot accuracy scores are provided in Table~\ref{tab:2t}; for evaluations involving both accuracy and normalized accuracy metrics, we present their mean value. The performance metrics for StableLM 3B at 2T tokens are sourced directly from its original report. Our results demonstrate that \text{BitNet b1.58} achieves top performance across all evaluated tasks, highlighting the strong generalization of 1.58-bit LLMs.

\section{Discussion and Future Work}

\textbf{1-bit Mixture-of-Experts (MoE) LLMs}

Mixture-of-Experts (MoE) architectures have established themselves as an efficient solution for scaling LLMs in a cost-effective manner. Despite reducing the computation FLOPs, MoE models still face substantial barriers due to large memory demands and significant inter-chip communication requirements. These limitations may be mitigated by leveraging 1.58-bit LLMs. By reducing memory utilization, fewer hardware resources are needed to run MoE models. Additionally, the transfer of activations between devices becomes far less burdensome. In the ideal case, when an entire model fits within a single chip, communication overhead would be essentially eliminated.

\textbf{Native Support of Long Sequence in LLMs}

Long sequence processing has emerged as an essential requirement for modern LLMs. A key obstacle for scaling sequence length during inference is the substantial memory required for storing KV caches. \text{BitNet b1.58} marks considerable progress in providing native support for lengthy input sequences by compressing activation precision from 16 bits to 8 bits. This reduction effectively permits the use of sequences twice as long given fixed hardware resources. Furthermore, lossless compression down to 4 bits or potentially less for 1.58-bit LLMs is feasible, which we propose as a direction for future exploration.

\textbf{LLMs on Edge and Mobile}

Adopting 1.58-bit LLMs holds great promise for significantly enhancing the deployment of language models on edge and mobile platforms. Constraints on memory capacity and processing power have historically hindered the execution of LLMs at the edge. The markedly lower memory and power demands of 1.58-bit LLMs, however, make them suitable for such devices, expanding the scope of feasible applications beyond current capabilities. This shift not only elevates the potential of edge and mobile technologies but also makes possible new practical use cases for LLMs. Importantly, 1.58-bit LLMs are inherently more compatible with CPUs—the predominant processing units in edge and mobile devices—enabling efficient and effective execution of \text{BitNet b1.58} and similar models.

\textbf{New Hardware for 1-bit LLMs}

Recent developments such as those by Groq\footnote{\url{https://groq.com/}} have showcased the viability and impact of developing specialized hardware (such as LPUs) to support LLMs. Looking ahead, we advocate for dedicated research and engineering efforts focused on designing next-generation hardware and systems optimized specifically for 1-bit LLMs, in alignment with the computational advances brought by BitNet~\cite{bitnet}.